\newcommand{\DoNotLoadEpstopdf}{}
\documentclass[11pt]{article}
\usepackage[T1]{fontenc}
\usepackage{lmodern}
\usepackage[margin=1.05in]{geometry}
\usepackage{amsmath,amssymb,amsthm,mathtools}
\usepackage{microtype}
\usepackage{enumitem,longtable,array,booktabs}
\usepackage{xcolor}
\usepackage[colorlinks=true,linkcolor=blue!45!black,citecolor=blue!45!black,urlcolor=blue!45!black]{hyperref}
\usepackage[nameinlink,noabbrev,capitalise]{cleveref}
\hypersetup{pdftitle={A Conditional Reduction for Algebraic Kuga-Satake Correspondences},pdfauthor={OpenAI},pdfsubject={Conditional generic recovery and specialization for projective K3 surfaces}}
\newtheorem{theorem}{Theorem}[section]
\newtheorem{proposition}[theorem]{Proposition}
\newtheorem{lemma}[theorem]{Lemma}
\newtheorem{corollary}[theorem]{Corollary}
\theoremstyle{definition}
\newtheorem{definition}[theorem]{Definition}
\theoremstyle{remark}
\newtheorem{remark}[theorem]{Remark}
\numberwithin{equation}{section}
\setlist[enumerate]{label=\textup{(\roman*)},leftmargin=*}
\setlist[itemize]{leftmargin=*}
\setlength{\emergencystretch}{2em}
\title{A Conditional Reduction for Algebraic Kuga--Satake Correspondences}
\author{OpenAI}
\date{September 10, 2026}
\begin{document}
\maketitle
\begin{abstract}
For a polarized K3 surface whose primitive cohomology has full orthogonal Hodge group, one algebraic correspondence inducing a nonzero map from that cohomology to the second cohomology of an abelian variety suffices to recover the prescribed full Kuga--Satake correspondence. We give an equivalent condition using holomorphic one-forms on a generically finite surface cover. If this input holds very generally in a polarized component, the prescribed correspondence is algebraic throughout that component, for every choice of standard data on the transcendental part. The existence of the initial correspondence remains a hypothesis.
\end{abstract}
\tableofcontents
\newpage
\section{Introduction and statements}\label{sec:introduction}

The Kuga--Satake construction associates an abelian variety to a polarized Hodge structure of K3 type. Its original purpose was to place the second cohomology of a K3 surface inside a tensor construction on the first cohomology of an abelian variety \cite{KugaSatake}. The resulting embedding is a Hodge morphism. Its realization by an algebraic correspondence is a further assertion, usually called the algebraicity of the Kuga--Satake correspondence; see \cite[Chapter 4]{Huybrechts} and \cite{vanGeemen}. This assertion is a particular instance of the rational Hodge conjecture.

We prove two conditional reductions for this algebraicity problem. At a point where the primitive cohomology has full orthogonal Hodge group, an algebraic correspondence inducing any nonzero map from primitive cohomology to the second cohomology of an abelian variety suffices to recover the \emph{specified} map into the \emph{full} even-Clifford target. Very-general existence of such a map then implies algebraicity of the specified correspondence at every point of the same polarized component. These reductions do not establish the existence of the initial map and therefore do not resolve the unrestricted Kuga--Satake conjecture.

The proof uses three devices which may be useful in related algebraicity questions. Independent operations on the two weight-one factors recover every multiplicity, and rational descent handles nonsplit Clifford algebras. Proper Hilbert parameter spaces specialize the \emph{entire prescribed flat class}, including its rational coefficients, over an analytic period chart. Finally, signed square roots compare quadratic polarizations through operations on the abelian factors. This last step does not require algebraicity of a Hodge endomorphism of the K3 surface.

\subsection{Conventions and the exact target}

All varieties in the statements are over $\mathbb C$, and cohomology is rational Betti cohomology unless a different field is displayed. We write $\operatorname{CH}^r(-)_{\mathbb Q}$ for codimension-$r$ cycles modulo rational equivalence with rational coefficients. A subscript records dimension, so $\operatorname{CH}_a$ denotes cycles of dimension $a$. We write $\operatorname{NS}(S)$ for the N\'eron--Severi group, $\operatorname{Hg}$ for the Hodge group, $\operatorname{Hom}_{\mathrm{HS}}$ for rational Hodge maps, and $F^\bullet$ for the Hodge filtration. The Tate twist $(1)$ shifts Hodge bidegrees by $(-1,-1)$. On a Hodge component of type $(p,q)$ the Hodge circle acts as $e^{i(p-q)\theta}$. The joint Hodge group of several structures is the smallest rational algebraic group containing their simultaneous Hodge circle. Equal-weight rational maps are Hodge maps exactly when they intertwine this circle.

Let $(S,h)$ be a primitively amply polarized K3 surface, with $h^2=2d$. Put
\[
 V=h^\perp\subset H^2(S,\mathbb Q),\qquad
 T=T(S)=\operatorname{NS}(S)_{\mathbb Q}^{\perp}\subset V.
\]
Orthogonality here is for the cup form $q$. The space $V$ has dimension $21$ and signature $(2,19)$, with the positive two-plane determined by $H^{2,0}(S)$. The space $T$ has the same period line and dimension at least two.

For the final target, fix a rational quadratic polarization $Q$ of $T$ and all standard Kuga--Satake data: the full even Clifford construction, its parity identification by right multiplication, a rational weight-one polarization, and any chosen isogenous realization. Denote this abelian variety by $A_S$, and set $W_S=H^1(A_S,\mathbb Q)$. The standard map $T(1)\to\operatorname{End}(W_S)$ and the stipulated identification $W_S^\vee\simeq W_S(1)$ determine
\begin{equation}\label{eq:target}
 \kappa_S:T\longrightarrow W_S\otimes W_S
 \lhook\joinrel\longrightarrow H^2(A_S\times A_S,\mathbb Q).
\end{equation}
The last arrow is the ordered K\"unneth inclusion. To realize this map means to find
\begin{equation}\label{eq:cycle-target}
 \Gamma_S\in\operatorname{CH}^2(S\times A_S\times A_S)_{\mathbb Q},
 \qquad
 (p_{AA})_*\bigl(\operatorname{cl}(\Gamma_S)\cup p_S^*x\bigr)
 =\kappa_S(x)\quad(x\in T).
\end{equation}
Here $\operatorname{cl}$ is the rational Betti cycle class, and $p_S$ and $p_{AA}$ are the projections to $S$ and $A_S\times A_S$. Every normalization in \eqref{eq:target} is retained. The target is never replaced by a simple spin factor.

\subsection{The conditional theorems}

\begin{theorem}[Recovery at a point with full Hodge group]\label{main:generic}
Suppose $\operatorname{Hg}(V)=\operatorname{SO}(V,q)$. The following conditions are equivalent.
\begin{enumerate}
\item There are an abelian variety $B$ and a cycle $\alpha\in\operatorname{CH}^2(S\times B)_{\mathbb Q}$ such that
\[
 \alpha_*:V\longrightarrow H^2(B,\mathbb Q)
\]
is nonzero.
\item There are a smooth connected projective surface $Y$ and a dominant generically finite morphism $f:Y\to S$ such that
\begin{equation}\label{eq:wedge-condition}
 f_*\!\left(\operatorname{im}\!\left[
 \bigwedge\nolimits^2 H^0(Y,\Omega_Y^1)
 \longrightarrow H^0(Y,\Omega_Y^2)\right]\right)\ne0.
\end{equation}
\item The prescribed full Kuga--Satake map is realized as in \eqref{eq:cycle-target}, for any stipulated standard data on $T$.
\end{enumerate}
In this situation $T=V$. The same input realizes the standard map for the full primitive Clifford construction with the cup form.
\end{theorem}

In \eqref{eq:wedge-condition}, $f_*$ is the trace on holomorphic two-forms, equivalently the restriction of cohomological pushforward to $H^{2,0}$. There is no genericity condition on $B$. The hypothesis on $V$ holds outside a countable union of proper closed analytic subsets in every full polarized period chart; see Proposition~\ref{gen:genericity}.

\begin{theorem}[Specialization in a polarized component]\label{main:specialization}
Fix a component of the moduli space of primitively amply polarized K3 surfaces of degree $2d$. Assume that condition \textup{(i)} or condition \textup{(ii)} of Theorem~\ref{main:generic} holds very generally in that component. Then every $(S,h)$ in the component satisfies \eqref{eq:cycle-target}, with all the stipulated standard data on $T(S)$.

Here it suffices that, about each desired point, the input holds outside a countable union of proper closed analytic subsets of a full polarized period chart. In particular an input outside countably many proper algebraic subsets of the moduli component suffices. The witness cycles, abelian varieties, or surface covers need not vary in a family.
\end{theorem}

\begin{corollary}\label{main:all-degrees}
If the hypothesis of Theorem~\ref{main:specialization} holds in the required polarized components in every degree, then \eqref{eq:cycle-target} holds for every projective K3 surface.
\end{corollary}

\begin{proof}
Every projective K3 surface has an ample integral class. Divide it to primitivity in $H^2(S,\mathbb Z)$; the resulting integral $(1,1)$ class is the first Chern class of a line bundle, a positive power of which is ample. Hence it is itself ample, and its square is a positive even integer. Apply Theorem~\ref{main:specialization} in its component.
\end{proof}

An input in one polarized component yields the conclusion in that component. Constructing a very-general correspondence or surface cover with the stated property remains the external existence problem. The hypotheses concern actual algebraic cycles; absolute Hodge or motivated classes do not supply this input.

\subsection{Proof outline}

Section~\ref{sec:cycles} proves the equivalence of the two inputs and records the available cycle operations. Section~\ref{cliff:section} fixes the full Clifford target and tensor conventions. Section~\ref{gen:section} recovers that target from one nonzero input by identifying the spin constituents and applying independent multiplicity operations. Section~\ref{spec:section} specializes the entire prescribed class in a local projective family. Section~\ref{data:section} restricts to $T$ and compares quadratic polarizations; Section~\ref{sec:conclusion} assembles the two theorems. Appendix~\ref{app:dependencies} records the established inputs and their roles.

\section{Cycle operations and the two inputs}\label{sec:cycles}

A correspondence $Z\in\operatorname{CH}^a(P\times Q)_{\mathbb Q}$ between smooth projective varieties acts by
\begin{equation}\label{eq:degree-convention}
 Z_*x=(p_Q)_*\bigl(\operatorname{cl}(Z)\cup p_P^*x\bigr):
 H^j(P)\longrightarrow H^{j+2e}(Q)(e),\qquad e=a-\dim P.
\end{equation}
The twist makes this a map of Hodge structures of weight $j$. Composition through a variety of dimension $r$ adds codimensions and subtracts $r$; the cycle class of the composition acts as the composite \cite{Fulton}. In particular a degree-preserving correspondence from a surface to degree two has codimension two.

\begin{lemma}[Weight-one operations]\label{cyc:h1}
For complex abelian varieties $P,Q$, pullback identifies
\[
 \operatorname{Hom}(Q,P)\otimes\mathbb Q
 \simeq
 \operatorname{Hom}_{\mathrm{HS}}\bigl(H^1(P,\mathbb Q),H^1(Q,\mathbb Q)\bigr).
\]
Thus arbitrary rational weight-one Hodge maps can act algebraically and independently on the two factors of an ordered tensor product of first cohomology groups.
\end{lemma}

\begin{proof}
Dualize a rational Hodge map to a map on first homology. After multiplication by a positive integer it preserves the integral lattices. The Hodge condition says that its real extension commutes with the complex structures on the universal covers. It therefore induces a holomorphic group homomorphism $Q\to P$, which is algebraic by projectivity. Conversely pullback by every such homomorphism is a Hodge map. This is the usual equivalence between abelian varieties up to isogeny and polarizable rational Hodge structures of weight one, written contravariantly \cite[Theorem 6.8 and Corollary 6.9]{Milne}. Transpose graphs realize the pullbacks. Products of two such correspondences, and their rational linear combinations, act independently on the two tensor factors. Preservation of the chosen polarizations is not required.
\end{proof}

\begin{lemma}[Degree projectors]\label{cyc:projectors}
Every cohomological K\"unneth projector on a product of abelian varieties is induced by a rational algebraic correspondence. If $U=H^1(B,\mathbb Q)$ and $m:B\times B\to B$ is addition, then the mixed part of $m^*$ gives an injection
\[
 H^2(B,\mathbb Q)=\bigwedge\nolimits^2U\longrightarrow U\otimes U,
 \qquad u\wedge v\longmapsto u\otimes v-v\otimes u.
\]
\end{lemma}

\begin{proof}
For an abelian variety $B$ of dimension $b$, multiplication by $2$ acts on $H^j(B)$ by $2^j$. Thus the rational polynomial
\[
 e_j(T)=\prod_{\substack{0\le k\le 2b\\k\ne j}}
 \frac{T-2^k}{2^j-2^k}
\]
applied to the correspondence $[2]^*$ is the projector onto $H^j(B)$. Apply these polynomials separately to each factor of a product. Finally $m^*u=p_1^*u+p_2^*u$ on $H^1$; multiplication in the exterior algebra, with its usual signs, gives the displayed formula.
\end{proof}

\begin{lemma}[Rational descent of a cycle span]\label{cyc:descent}
Let $E$ be a finite-dimensional rational vector space of cohomological maps, and let $E_{\mathrm{alg}}\subset E$ be the image of rational algebraic correspondences of a fixed degree. For every field extension $K/\mathbb Q$,
\[
 E\cap(E_{\mathrm{alg}}\otimes_{\mathbb Q}K)=E_{\mathrm{alg}}
 \quad\text{inside }E\otimes_{\mathbb Q}K.
\]
The same assertion holds if only a specified rational subspace of the source is considered.
\end{lemma}

\begin{proof}
The cycle class map and restriction of its action are rational linear maps, so $E_{\mathrm{alg}}$ is a rational subspace. If a rational target lies in the scalar-extended span of finitely many rational cycle actions, write these actions as columns in a rational basis of $E$. The linear system with that target as right-hand side is consistent over $K$, and row reduction over $\mathbb Q$ gives a rational solution. Only finitely many cycles are used.
\end{proof}

\begin{proposition}[Correspondences and surface covers]\label{cyc:equivalence}
Conditions \textup{(i)} and \textup{(ii)} of Theorem~\ref{main:generic} are equivalent. The implication from the surface condition to a nonzero correspondence on $V$ does not require the full Hodge-group hypothesis.
\end{proposition}

\begin{proof}
Suppose first that $f:Y\to S$ satisfies \eqref{eq:wedge-condition}. Let $D=\operatorname{Alb}(Y)$, let $a:Y\to D$ be an Albanese map, and put $b=\dim D$. Albanese pullback identifies holomorphic one-forms \cite[Chapter VI, \S9.2, pp.~315--316]{Albanese}. Hence there is
\[
 \beta\in H^{2,0}(D)=\bigwedge\nolimits^2H^{1,0}(D),\qquad
 y=f_*a^*\beta\ne0.
\]
In particular $b\ge2$ and $y\in H^{2,0}(S)\subset V_{\mathbb C}$. The surface correspondence $(a,f)_*[Y]$ has codimension $b$ on $D\times S$. Its transpose acts by
\[
 a_*f^*:H^2(S)\longrightarrow H^{2b-2}(D)(b-2).
\]
This action is nonzero on $V$: the projection formula gives
\begin{equation}\label{eq:albanese-nonzero}
 \int_D\beta\cup a_*f^*\overline y
 =\int_Ya^*\beta\cup f^*\overline y
 =\int_Sy\cup\overline y\ne0.
\end{equation}

We give an algebraic degree adjustment. For an ample line bundle $L$ on $D$, written additively, set
\[
 E=m^*L-p_1^*L-p_2^*L\in\operatorname{CH}^1(D\times D),
\]
where $m$ is addition. The class of $E$ is purely mixed. Put $U=H^1(D,\mathbb Q)$ and write this class as $\sum_i u_i\otimes v_i$, choosing a basis $(u_i)$ of $U$. Its matrix
\[
 M:U^\vee\longrightarrow U,\qquad
 M(\varphi)=\sum_i\varphi(u_i)v_i
\]
is invertible: it is the alternating tensor associated with $c_1(L)$, whose top exterior power is nonzero since $L$ is ample. Intrinsically it is a Hodge map $M:U^\vee\to U(1)$; in the matrix calculation we use the underlying rational spaces. With the Koszul convention for K\"unneth products,
\[
 \operatorname{cl}(E)^2
 =-2\sum_{i<j}(u_i\wedge u_j)\otimes(v_i\wedge v_j).
\]
Poincar\'e duality identifies $H^{2b-2}(D)$ with $(\bigwedge^2U)^\vee$ on underlying rational spaces, by integration against degree-two classes. Consequently $(E^2)_*$ is $-2\bigwedge^2M$ after this identification, and is an isomorphism
\begin{equation}\label{eq:degree-adjustment}
 (E^2)_*:H^{2b-2}(D)(b-2)\xrightarrow{\ \sim\ }H^2(D).
\end{equation}
The twists can also be checked from
$H^{2b-2}(D)(b-2)\simeq H^2(D)^\vee(-2)$ and
$\bigwedge^2M:H^2(D)^\vee\to H^2(D)(2)$.
Composing with $(f,a)_*[Y]$ yields a cycle of codimension $b+2-b=2$ on $S\times D$, nonzero on $V$ by \eqref{eq:albanese-nonzero}. This uses a square of an explicit divisor; it assumes no algebraic inverse Lefschetz operator for a general variety.

Conversely suppose that $\alpha_*:V\to H^2(B)$ is nonzero and $\operatorname{Hg}(V)=\operatorname{SO}(V,q)$. The standard representation is irreducible, so $\alpha_*|_V$ is injective. For $0\ne\sigma\in H^{2,0}(S)$, put $\eta=\alpha_*\sigma\ne0$ in $H^{2,0}(B)$. Write $c=\dim B\ge2$ and let $\ell$ be an ample divisor class on $B$, replaced by a very ample multiple if necessary. The rational surface cycle
\[
 C=\alpha\cdot p_B^*\ell^{c-2}\in\operatorname{CH}_2(S\times B)_{\mathbb Q}
\]
satisfies
\begin{equation}\label{eq:surface-nonzero}
 \int_{S\times B}\operatorname{cl}(C)\cup p_S^*\overline\sigma\cup p_B^*\eta
 =\int_B\overline\eta\cup\eta\cup\ell^{c-2}\ne0.
\end{equation}
Here rationality gives $\alpha_*\overline\sigma=\overline\eta$, and the last pairing is definite up to a fixed nonzero sign on $H^{2,0}(B)$ by the Hodge--Riemann bilinear relations. An actual representative of $C$ is a finite rational combination of integral surfaces; for example use general hyperplanes to represent the indicated intersections. At least one component has nonzero pairing in \eqref{eq:surface-nonzero}. Resolve that component projectively and call the resulting smooth connected surface $Y$. Its projections give morphisms $f:Y\to S$ and $a':Y\to B$ with
\[
 \int_Y f^*\overline\sigma\cup a'^*\eta\ne0.
\]
Thus $f^*\sigma\ne0$, which forces $f$ to have two-dimensional image. It is therefore dominant and generically finite. The projection formula also gives $f_*a'^*\eta\ne0$. Since $\eta$ is a sum of wedges of holomorphic one-forms on $B$, its pullback lies in the wedge image in \eqref{eq:wedge-condition}. This proves the converse.
\end{proof}

\begin{remark}
The reverse implication uses injectivity on $V$, not merely nonvanishing on some divisor class. At a Picard jump a correspondence nonzero only on new algebraic classes need not satisfy the surface condition. No such assertion is used below.
\end{remark}

\section{Full Clifford structures and their tensor maps}\label{cliff:section}

We fix the conventions needed to identify the final correspondence exactly.
The classical construction is reviewed in
\cite[Chapter~4, \S2]{Huybrechts} and \cite[\S\S5--6]{vanGeemen}; we
include the details concerning signs, polarizations, and the two tensor slots.
Our Hodge circle acts on a vector of type $(p,q)$ by
$e^{i(p-q)\theta}$. Tate twists leave this circle unchanged.

Let $H$ be a rational Hodge structure of K3 type: it has weight two,
$h^{2,0}=h^{0,2}=1$, and no Hodge types other than $(2,0)$, $(1,1)$,
and $(0,2)$. Let $Q$ be a
nondegenerate rational quadratic Hodge form. We allow either overall sign:
$Q$ has sign $\epsilon\in\{1,-1\}$ on the real period plane and the
opposite sign on its orthogonal. Write
\[
 C(H,Q)=C^+(H,Q)\oplus C^-(H,Q),\qquad v^2=Q(v,v).
\]
If $H^{2,0}=\mathbb C(e_1+i e_2)$, where $e_1,e_2$ are real orthogonal
vectors with $Q(e_j,e_j)=\epsilon$, set
\begin{equation}\label{cliff:circle}
 J=\epsilon e_1e_2,\qquad
 g_\theta=\cos\theta+J\sin\theta.
\end{equation}
The element $J$ depends only on the oriented period plane. Left
multiplication by $g_\theta$ gives both parity spaces weight-one Hodge
structures, denoted $P_H^+$ and $P_H^-$, with the $i$-eigenspace of $L_J$
as their $(1,0)$-part. We retain these \emph{full} spaces throughout.

\begin{proposition}\label{cliff:construction}
The Hodge structures $P_H^\pm$ are polarizable, and
\begin{equation}\label{cliff:multiplication}
 m^Q:H(1)\longrightarrow\operatorname{Hom}(P_H^+,P_H^-),
 \qquad x\longmapsto L_x
\end{equation}
is a rational Hodge embedding. The parity spaces are rationally
Hodge-isomorphic by right multiplication by any fixed invertible odd
Clifford element. On a sufficiently small period chart, their Hodge
filtrations vary holomorphically and admit fixed rational polarizations.
\end{proposition}

\begin{proof}
The Clifford relations give
\[
 J^2=-1,\quad [J,e_1]=-2e_2,\quad [J,e_2]=2e_1,
 \quad [J,e_1+i e_2]=2i(e_1+i e_2).
\]
Moreover, $J$ commutes with the orthogonal complement of the period
plane. Thus $g_\theta\in\operatorname{Spin}(H,Q)(\mathbb R)$, and its
conjugation action on $H$ is the prescribed weight-two Hodge circle.
Since
\[
 L_{g_\theta}L_xL_{g_\theta}^{-1}
   =L_{g_\theta xg_\theta^{-1}},
\]
the map in \eqref{cliff:multiplication} is a Hodge map. It is injective
because $L_x(1)=x$.
Choose a rational nonisotropic vector $z\in H$. Right multiplication
$R_z$ exchanges the parities, has inverse $R_{z/Q(z,z)}$, and commutes
with the left circle. The same observation applies to any invertible
odd right multiplier.

For holomorphic variation, put $\omega=e_1+i e_2$.
The relation $J\omega=i\omega$ shows that the image of $L_\omega$
from the opposite parity is contained in the $i$-eigenspace of $L_J$.
Conversely $\omega J=-i\omega$, so this eigenspace is contained in
$\ker L_\omega$. Finally
\[
 L_\omega L_{\bar\omega}+L_{\bar\omega}L_\omega
       =4\epsilon\operatorname{id}
\]
shows that $\ker L_\omega=\operatorname{im}L_\omega$, with the
appropriate parities understood. Hence
$F^1P_{H,\mathbb C}^\pm=\operatorname{im}L_\omega$; this is a
holomorphic subbundle as the isotropic line varies.

Here is a fixed rational polarization. Write $u\mapsto u^*$ for
reversal in the Clifford algebra and $\operatorname{scal}$ for its
scalar part. Choose rational orthogonal vectors $f_1,f_2$ of sign
$\epsilon$ whose oriented plane is sufficiently close to the given
period plane. Rational approximation followed by rational
orthogonalization supplies them. Put $a=f_1f_2$ and, on $C^+(H,Q)$,
\begin{equation}\label{cliff:polarizing-form}
 s(u,v)=\operatorname{scal}(uav^*).
\end{equation}
The scalar part is cyclic and invariant under reversal, as is checked
on an orthogonal monomial basis. Since $a^*=-a$, the form $s$ is
alternating. It is nondegenerate because $a$ is invertible and
$(u,v)\mapsto\operatorname{scal}(uv^*)$ is nondegenerate. Reversal
gives $g_\theta^*g_\theta=1$, so $s$ is invariant under the left circle
at every period.

At the reference plane spanned by $f_1,f_2$, write $a=cJ$ with
$c\in\mathbb R^\times$. Then
\[
 s(u,Jv)=c\operatorname{scal}\bigl(\operatorname{Ad}(J)u\,v^*\bigr).
\]
To see definiteness, extend a normalized basis of this plane to an
orthogonal basis with squared lengths
$\epsilon,\epsilon,-\epsilon,\ldots,-\epsilon$.
The last pairing is diagonal on even monomials $e_I$.
If $r=|I|$ and $k$ entries of $I$ belong to the period plane, its
diagonal sign is
\[
 (-1)^k\epsilon^r(-1)^{r-k}=1,
\]
because $r$ is even. Up to the overall sign required by the
weight-one polarization convention, $s$ is therefore a polarization.
Circle invariance makes $s(u,Jv)$ symmetric and nondegenerate at every
period. Its signature is constant on a connected period chart
containing both the reference plane and the original period, so the
chosen sign of $s$ polarizes throughout that chart.
Transport by $R_z$ supplies a fixed polarization
on $C^-(H,Q)$ as well.
\end{proof}

Choose abelian varieties realizing the polarizable weight-one spaces
in Proposition~\ref{cliff:construction}; rational realizations suffice.
For a weight-one Hodge structure $P$, fix the rational Hodge
isomorphism
\[
 j_P:P^\vee(-1)\xrightarrow{\sim}P
\]
induced by its chosen polarization. A map
$f:H(1)\to\operatorname{Hom}(P,M)$ then determines the weight-two
tensor map
\begin{equation}\label{cliff:tensor}
 H\longrightarrow\operatorname{Hom}(P,M)(-1)
   =M\otimes P^\vee(-1)
   \xrightarrow{1\otimes j_P}M\otimes P.
\end{equation}
When $H\subset H^2(S,\mathbb Q)$ for a surface $S$, choose abelian
realizations $M=H^1(A_M,\mathbb Q)$ and $P=H^1(A_P,\mathbb Q)$.
We call $f$ \emph{algebraically realized} if some cycle in
$\operatorname{CH}^2(S\times A_M\times A_P)_{\mathbb Q}$ has, on
$H$, exactly the action \eqref{cliff:tensor} included in the
$H^1\otimes H^1$ K\"unneth component. This specifies the restricted
action on $H$; the algebraic target projectors may be applied to
remove other target components. The terminology always refers to
actual algebraic cycles.

\begin{lemma}[Operations on the two slots]\label{cliff:slots}
Suppose $f:H(1)\to\operatorname{Hom}(P,M)$ is algebraically realized.
For rational weight-one Hodge maps
$a:M\to M'$ and $b:P'\to P$, the map
$x\mapsto a f(x)b$ is algebraically realized. The same holds after
changing the chosen duality isomorphisms, interchanging the tensor
factors, or making rational weight-one Hodge identifications with
isogenous abelian models.
\end{lemma}

\begin{proof}
Under \eqref{cliff:tensor}, the two operations are precisely
\begin{equation}\label{cliff:slot-formula}
 a\otimes\bigl(j_{P'}b^\vee(-1)j_P^{-1}\bigr).
\end{equation}
Both factors are rational weight-one Hodge maps, so
Lemma~\ref{cyc:h1} realizes them by abelian correspondences.
Their product acts independently on the two factors and preserves
the codimension of the correspondence from the surface.
Changing $j_P$ likewise composes with a weight-one Hodge
isomorphism. The factor-switching isomorphism is algebraic; its
K\"unneth sign can be included in the rational coefficient.
The projectors of Lemma~\ref{cyc:projectors} remove all other
K\"unneth components when necessary.
\end{proof}

In particular, the standard full even-Clifford model is $W=P_H^+$.
For a stipulated invertible odd right multiplier $z$, its
endomorphism-valued embedding is exactly
\begin{equation}\label{cliff:standard}
 H(1)\longrightarrow\operatorname{End}(W),\qquad
 x\longmapsto\bigl[w\longmapsto xwz\bigr]
       =R_zL_x.
\end{equation}
Applying the stipulated duality $W^\vee\simeq W(1)$ gives its
specified map to $W\otimes W$. Thus algebraicity of $m^Q$ gives
algebraicity of this precise map, with the fixed multiplier,
polarization, rational normalization, and identification with
$H^1$ of the chosen Kuga--Satake abelian variety. No simple factor
replaces the full even-Clifford space.

These conclusions are also independent of the usual equivalent
Clifford conventions. Reversal takes left multiplication by a spin
element to right multiplication by its inverse, since
$(gu)^*=u^*g^{-1}$. Passing between homology and cohomology uses
duality, transpose, and a weight-one polarization. These are exactly
the operations in Lemma~\ref{cliff:slots}; the source Hodge circle
is always the one fixed above. Reversing the period Hodge action
independently would be a different construction.

\section{Recovering the full Clifford correspondence at a general point}
\label{gen:section}

The argument in this section uses the equality
\begin{equation}\label{gen:full-group}
 \operatorname{Hg}(V)=\operatorname{SO}(V,q),
 \qquad V=h^\perp\subset H^2(S,\mathbb Q),\qquad \dim V=21.
\end{equation}
There is no genericity assumption on the abelian variety that supplies the
initial algebraic correspondence.  The principal point is to recover every
multiplicity in the full Clifford construction from one nonzero map.

\subsection{The generic Hodge group}

\begin{proposition}\label{gen:genericity}
In any sufficiently small, full-dimensional marked polarized period chart,
Equation~\eqref{gen:full-group} holds outside a countable union of proper
closed analytic subsets.  At every point satisfying
Equation~\eqref{gen:full-group}, the Hodge structure $V$ is irreducible and
$T(S)=V$.
\end{proposition}

\begin{proof}
Polarized local Torelli identifies such a chart with an open subset of the
domain of oriented positive two-planes in the fixed real quadratic space
$V_{\mathbb R}$; see \cite[Chapter~6]{Huybrechts}.
The infinitesimal generators of the Hodge circles over any such open
subset span $\mathfrak{so}(V_{\mathbb R},q)$.  Indeed, under the natural
identification $\mathfrak{so}(V_{\mathbb R},q)\simeq\Lambda^2V_{\mathbb R}$,
each generator is a nonzero multiple of the bivector of its positive
plane.  A linear functional vanishing on all these bivectors is an
alternating bilinear form vanishing on an open set of ordered pairs of
vectors.  It therefore vanishes identically.

The Hodge group is a connected rational algebraic subgroup of
$\operatorname{SO}(V,q)$.  Connectedness follows because its rational
identity component contains the connected Hodge circle and hence, by
minimality, is the whole group.  If this subgroup is proper, its Lie
algebra is a nonzero proper rational subspace
$L\subset\mathfrak{so}(V,q)$.  Set $r=\dim L$, and choose a rational
generator $\tau_L$ of $\Lambda^rL$.  Conjugation by the circle preserves
$L$ and acts trivially on its real determinant line: a continuous
character from a compact connected circle to $\mathbb R^\times$ is
trivial.  Thus $\tau_L$, viewed in
$\Lambda^r\operatorname{End}(V)$, has Hodge type $(0,0)$.

There are only countably many possible rational subspaces $L$.
For each one, the condition that $\tau_L$ have type $(0,0)$ is closed
analytic in the period chart.  In fact the tensor has weight zero and is
real, so this condition is equivalent to $\tau_L\in F^0$, a holomorphic
condition.  The locus is proper.  Otherwise $L_{\mathbb R}$ would be
preserved by every circle in the chart, hence by the adjoint actions of
their infinitesimal generators.  The spanning assertion above would
make $L_{\mathbb C}$ an ideal of the simple Lie algebra
$\mathfrak{so}_{21}(\mathbb C)$, a contradiction.

Finally, the standard representation of $\operatorname{SO}(V,q)$ is
absolutely irreducible and has no nonzero invariants.  Consequently $V$
has no proper nonzero rational Hodge substructure and no nonzero
rational vector of type $(1,1)$.  Since the primitive part of
$\operatorname{NS}(S)_{\mathbb Q}$ consists of such vectors,
$\operatorname{NS}(S)_{\mathbb Q}=\mathbb Qh$ and $T(S)=V$.
\end{proof}

\subsection{The spin representation forced by weight one}

Assume Equation~\eqref{gen:full-group}.  Let $B$ be a complex abelian
variety, put $U=H^1(B,\mathbb Q)$, and write
\[
 W^+=C^+(V,q),\qquad W^-=C^-(V,q)
\]
for the full Clifford parity spaces with their weight-one Hodge
structures.  The Hodge circle on these spaces acts by left multiplication
by its lift to $\operatorname{Spin}(V,q)$.
Let $\mathcal H$ be the joint Hodge group of $V,U,W^+,W^-$, and let
$\mathfrak h$ be its rational Lie algebra.

\begin{lemma}\label{gen:joint-group}
There is a decomposition into commuting rational ideals
\[
 \mathfrak h=\mathfrak g\oplus\mathfrak k,
 \qquad \mathfrak g\simeq\mathfrak{so}(V,q),
\]
where $\mathfrak k$ is the kernel of the action on $V$.
The ideal $\mathfrak k$ acts trivially on $W^+$ and $W^-$, and
$\mathfrak g$ acts on them by the usual left spin action.
\end{lemma}

\begin{proof}
The group $\mathcal H$ is connected by the same circle argument as in
Proposition~\ref{gen:genericity}.  It is reductive.  To see this directly,
each of the four pure summands is semisimple as a rational
$\mathcal H$-module: an invariant rational subspace is a Hodge
substructure, and a polarization provides a rational Hodge complement.
This is the usual semisimplicity argument for polarizable Hodge
structures; compare \cite[Proposition~3.6]{Deligne}.  Their direct sum is
a faithful semisimple module.  The unipotent radical has nonzero
invariants in each irreducible constituent; normality makes those
invariants $\mathcal H$-stable, so it acts trivially on every
constituent.  Faithfulness forces the radical to be trivial.

Projection to $V$ maps $\mathcal H$ onto
$\operatorname{Hg}(V)=\operatorname{SO}(V,q)$.
Indeed, its closed rational image contains the projected circle and
hence contains $\operatorname{Hg}(V)$. Conversely, the preimage of
$\operatorname{Hg}(V)$ contains the joint circle, so minimality places
$\mathcal H$ inside that preimage.
The reductive rational Lie algebra $\mathfrak h$ is
the direct sum of its center and its rational simple ideals.  Exactly
one simple ideal maps nontrivially to the absolutely simple quotient
$\mathfrak{so}(V,q)$: its image is the whole quotient, and the commuting
images of all other ideals must vanish.  The nonzero map from that
simple ideal is an isomorphism.  This gives the asserted rational
decomposition.

On $V\oplus W^+\oplus W^-$, the image of $\mathcal H$ is contained in
the rational algebraic image of $\operatorname{Spin}(V,q)$ acting on
$V$ through its orthogonal quotient and on $W^\pm$ by left
multiplication.  That image contains the projected Hodge circle, so
the containment again follows by minimality.  The differential of
$\operatorname{Spin}(V,q)\to\operatorname{SO}(V,q)$ is an isomorphism.
Its kernel on Lie algebras is therefore zero, proving both remaining
assertions.
\end{proof}

Let $\Delta$ denote the complex spin representation of
$\mathfrak{so}_{21}(\mathbb C)$, of dimension $2^{10}$.

\begin{lemma}\label{gen:width}
There are complex multiplicity spaces $M,N_+,N_-$ and a trivial
$\mathfrak g_{\mathbb C}$-module $M_0$ such that
\begin{equation}\label{gen:spin-decompositions}
 U_{\mathbb C}=(\Delta\otimes M)\oplus M_0,
 \qquad W^\pm_{\mathbb C}=\Delta\otimes N_\pm,
 \qquad \dim N_\pm=2^{10}.
\end{equation}
Every rational $\mathfrak g$-equivariant map $U\to W^\pm$ is a
weight-one Hodge map.  Consequently
\begin{equation}\label{gen:hom-base-change}
 \operatorname{Hom}_{\mathrm{HS}}(U,W^\pm)\otimes_{\mathbb Q}\mathbb C
 =\operatorname{Hom}_{\mathfrak g_{\mathbb C}}
       (U_{\mathbb C},W^\pm_{\mathbb C})
 \simeq\operatorname{Hom}_{\mathbb C}(M,N_\pm).
\end{equation}
\end{lemma}

\begin{proof}
Let $D$ be the infinitesimal Hodge-circle generator divided by $i$;
its eigenvalue on type $(p,r)$ is $p-r$.  In the decomposition of
Lemma~\ref{gen:joint-group}, write $D=D_{\mathfrak g}+D'$.  On
$V_{\mathbb C}$ the eigenvalues are $2,-2,0,\ldots,0$.
In the standard Cartan coordinates for type $B_{10}$, in which the
vector weights are $\pm e_1,\ldots,\pm e_{10},0$, we may therefore take
$D_{\mathfrak g}=2e_1^\vee$.

The element $D_{\mathfrak g}$ is semisimple on every finite-dimensional
$\mathfrak g_{\mathbb C}$-module.  It commutes with $D$, so
$D'=D-D_{\mathfrak g}$ is also semisimple on $U_{\mathbb C}$.
Each $D'$-eigenspace is $\mathfrak g_{\mathbb C}$-stable.
Consider an irreducible constituent in such an eigenspace, with
highest weight $\lambda=(\lambda_1,\ldots,\lambda_{10})$.
The highest-weight description for $B_{10}$ gives
\[
 \lambda_1\ge\cdots\ge\lambda_{10}\ge0,
 \qquad
 \lambda_i\in\mathbb Z\ \text{for every }i,
 \quad\text{or}\quad
 \lambda_i\in\tfrac12+\mathbb Z\ \text{for every }i.
\]
The largest and smallest eigenvalues of $2e_1^\vee$ are respectively
$2\lambda_1$ and $-2\lambda_1$.  Indeed $e_1^\vee$ is dominant, and the
Weyl group permits every coordinate sign change.  These elementary
highest-weight and spin facts may be found in the orthogonal-group
chapters of \cite{FultonHarris}.

Adding the scalar by which $D'$ acts does not change the difference
$4\lambda_1$ between the extreme eigenvalues.  Since $U$ has Hodge
types $(1,0)$ and $(0,1)$, the eigenvalues of $D$ on $U_{\mathbb C}$
belong to $\{1,-1\}$.  Hence $4\lambda_1\le2$.  The only possibilities
are the trivial representation and
$\lambda=(\tfrac12,\ldots,\tfrac12)$, which is $\Delta$.
On a spin constituent the extreme eigenvalues already are $1,-1$;
the only scalar shift preserving this pair is zero.  Thus
\begin{equation}\label{gen:kernel-circle-zero}
 D'=0\quad\text{on }\Delta\otimes M.
\end{equation}
This proves the assertion about $U$.  The complex Clifford identity
$C^+(V_{\mathbb C},q)\simeq\operatorname{End}_{\mathbb C}(\Delta)$,
with left multiplication, gives the assertion about $W^+$.
An invertible odd right multiplier identifies the two parity spaces
for the left spin action and gives the assertion about $W^-$.
Lemma~\ref{gen:joint-group} shows that $D'=0$ on both of them.

Now a rational $\mathfrak g$-map $f:U\to W^\pm$ vanishes on $M_0$
after complexification.  The summand $M_0$ is $D$-stable.  On the spin
summand, Equation~\eqref{gen:kernel-circle-zero} shows that $f$
intertwines $D$.  It therefore intertwines $D$ on all of $U$, hence
the Hodge circle, and is a Hodge map.  Conversely, every rational
Hodge map intertwines $\mathcal H$ and therefore $\mathfrak g$.

Finally, rational $\mathfrak g$-intertwiners are the kernel of the
finite system of rational linear equations
\[
 f\rho_U(X_a)=\rho_{W^\pm}(X_a)f,
\]
where $\{X_a\}$ is a rational basis of $\mathfrak g$.
Taking that kernel commutes with extension of scalars.  Schur's
Lemma over $\mathbb C$ and Equation~\eqref{gen:spin-decompositions}
now give Equation~\eqref{gen:hom-base-change}.
\end{proof}

\begin{lemma}\label{gen:vector-multiplicity}
The space
$\operatorname{Hom}_{\mathfrak g_{\mathbb C}}
 (V_{\mathbb C},\Delta\otimes\Delta)$ has dimension one.
Neither $\Delta$ nor the trivial representation contains a vector
copy $V_{\mathbb C}$.
\end{lemma}

\begin{proof}
The spin representation in type $B_{10}$ is self-dual, since the
longest Weyl element is $-1$.  Under the conjugation action, the
Clifford antisymmetrization map gives
\begin{equation}\label{gen:clifford-exterior}
 \Delta\otimes\Delta\simeq\operatorname{End}(\Delta)
 \simeq C^+(V_{\mathbb C},q)
 \simeq\bigoplus_{k\ \mathrm{even}}\Lambda^kV_{\mathbb C}.
\end{equation}
Here the action is conjugation, in contrast to the left action used
in Equation~\eqref{gen:spin-decompositions}.

For completeness, the vector multiplicity in
Equation~\eqref{gen:clifford-exterior} can be read off without
decomposing every exterior power.  Choose an orthonormal basis
$v_1,\ldots,v_{21}$.  An equivariant map
$F:V_{\mathbb C}\to\Lambda^kV_{\mathbb C}$ is determined by $F(v_1)$,
because the orthogonal orbit of $v_1$ spans $V_{\mathbb C}$.
Its value is fixed by all determinant-one diagonal sign changes on
$v_2,\ldots,v_{21}$.  A basis exterior monomial fixed by all these
sign changes contains either none or all of these twenty vectors.
For even $k$, only $k=0$ and $k=20$ can therefore occur.
The degree-zero map vanishes, since a determinant-one orthogonal
transformation sending $v_1$ to $-v_1$ forces
$F(v_1)=F(-v_1)=-F(v_1)$.  For $k=20$ the possible value is a multiple
of $v_2\wedge\cdots\wedge v_{21}$, and volume duality
$\Lambda^{20}V_{\mathbb C}\simeq V_{\mathbb C}$ supplies a nonzero
equivariant map.  The multiplicity is exactly one.
The last assertion follows from irreducibility and the distinct
dimensions $21$, $2^{10}$, and $1$.
\end{proof}

\subsection{Independent target operations and rational descent}

\begin{theorem}\label{gen:generation}
Assume Equation~\eqref{gen:full-group}.  Suppose that
$\alpha\in\operatorname{CH}^2(S\times B)_{\mathbb Q}$ induces a
nonzero map $V\to H^2(B,\mathbb Q)$ for some complex abelian variety
$B$.  Let $\epsilon_1,\epsilon_2\in\{+,-\}$, and let $A_1,A_2$ be
abelian varieties realizing $W^{\epsilon_1},W^{\epsilon_2}$ as their
first rational cohomology.  Every rational Hodge map
\[
 \beta:V\longrightarrow W^{\epsilon_1}\otimes W^{\epsilon_2}
 \subset H^2(A_1\times A_2,\mathbb Q)
\]
is the restriction of the action of a cycle
$\Gamma\in\operatorname{CH}^2(S\times A_1\times A_2)_{\mathbb Q}$.
In particular, this recovers the full primitive Clifford
multiplication tensor and its prescribed even-Clifford realization.
\end{theorem}

\begin{proof}
Pull back $\alpha$ along addition $B\times B\to B$ and apply the
algebraic projector to $U\otimes U\subset H^2(B\times B,\mathbb Q)$
from Lemma~\ref{cyc:projectors}.  The resulting codimension-two cycle
induces
\[
 \alpha':V\longrightarrow U\otimes U.
\]
It is nonzero: on $H^2(B,\mathbb Q)=\Lambda^2U$, the mixed component
of addition pullback is the injective map
$u\wedge v\mapsto u\otimes v-v\otimes u$.

Fix a nonzero $\mathfrak g_{\mathbb C}$-equivariant map
$\gamma:V_{\mathbb C}\to\Delta\otimes\Delta$.
The map $\alpha'$ is rational and Hodge, hence
$\mathfrak g$-equivariant.  Lemmas~\ref{gen:width} and
\ref{gen:vector-multiplicity} give
\begin{equation}\label{gen:initial-tensor}
 \alpha'_{\mathbb C}=\gamma\otimes t,
 \qquad 0\ne t\in M\otimes M.
\end{equation}
In this formula the tensor factors have been regrouped so that the
two spin factors come first.  The other summands of
$U_{\mathbb C}\otimes U_{\mathbb C}$ contain no vector copy.

Put $N_j=N_{\epsilon_j}$ for $j=1,2$.  The independent operations
\[
 a\in\operatorname{Hom}_{\mathbb C}(M,N_1),\qquad
 b\in\operatorname{Hom}_{\mathbb C}(M,N_2)
\]
send $t$ to $(a\otimes b)t$, and these images span
$N_1\otimes N_2$.  To check this, choose
$\mu,\nu\in M^\vee$ with $c=(\mu\otimes\nu)(t)\ne0$.
For arbitrary $n_1\in N_1$ and $n_2\in N_2$, the maps
\[
 a(m)=\mu(m)n_1,\qquad b(m)=\nu(m)n_2
\]
send $t$ to $c\,n_1\otimes n_2$.  Thus any symmetry imposed on $t$
by the original exterior square causes no restriction: the two
operations are chosen independently.

To implement these operations over $\mathbb Q$, choose rational bases
\[
 f_1,\ldots,f_r\quad\text{of }
 \operatorname{Hom}_{\mathfrak g}(U,W^{\epsilon_1}),\qquad
 g_1,\ldots,g_s\quad\text{of }
 \operatorname{Hom}_{\mathfrak g}(U,W^{\epsilon_2}).
\]
By Lemma~\ref{gen:width}, all these maps are Hodge, and their complex
spans are the full spaces of the independent operations just used.
Lemma~\ref{cyc:h1} realizes each by a rational abelian homomorphism
acting contravariantly on $H^1$; this is the classical weight-one
equivalence, also stated in \cite[Theorem~6.8 and Corollary~6.9]{Milne}.
Products of their transpose-graph correspondences, followed by the mixed
projector if necessary, therefore realize every map
\begin{equation}\label{gen:rational-generators}
 \Phi_{ij}=(f_i\otimes g_j)\alpha'
\end{equation}
by a rational algebraic cycle $\Gamma_{ij}$ on
$S\times A_1\times A_2$.
These cycles have codimension two: the target correspondences have
degree shift zero, and composition preserves the degree shift of
the original surface correspondence.

Lemma~\ref{gen:vector-multiplicity},
Equation~\eqref{gen:initial-tensor}, and the spanning calculation show
that the complex span of the maps in
Equation~\eqref{gen:rational-generators} is
\[
 \operatorname{Hom}_{\mathfrak g_{\mathbb C}}
 \bigl(V_{\mathbb C},
       W^{\epsilon_1}_{\mathbb C}\otimes
       W^{\epsilon_2}_{\mathbb C}\bigr).
\]
The prescribed rational Hodge map $\beta$ lies in this space.
Consequently the finite system
\[
 \beta=\sum_{i=1}^r\sum_{j=1}^s c_{ij}\Phi_{ij}
\]
has a solution over $\mathbb C$.  All its coefficients and its
right-hand side are rational, so it has a solution with
$c_{ij}\in\mathbb Q$ by Lemma~\ref{cyc:descent}.  The cycle
$\Gamma=\sum_{i,j}c_{ij}\Gamma_{ij}$ has precisely the required action.

Finally, Clifford multiplication is the rational Hodge map
\[
 V(1)\longrightarrow\operatorname{Hom}(W^+,W^-),
 \qquad x\longmapsto L_x.
\]
Using a weight-one polarization to identify
$(W^+)^\vee(-1)$ with $W^+$ makes it a rational Hodge map
$V\to W^-\otimes W^+$.  It is therefore one of the maps just realized.
The fixed odd right multiplier and prescribed weight-one duality
then give exactly the specified even-Clifford tensor, by
Lemma~\ref{cliff:slots}.  No choice of a single spin copy or
undetermined scalar remains.
\end{proof}

\begin{remark}\label{gen:rationality}
The complex decompositions in
Equation~\eqref{gen:spin-decompositions} are used only to compute a
span.  The maps $f_i,g_j$ and the cycles $\Gamma_{ij}$ are rational
before that computation.  In particular, the argument requires
neither a split rational Clifford algebra nor descent of an
individual complex spin summand.  It proves algebraicity of the
prescribed rational tensor on the full Clifford spaces.
\end{remark}

\section{Specialization of an actual algebraic class}
\label{spec:section}

We first establish a specialization statement for projective analytic
families.  This permits the Kuga--Satake construction to be made on a
marked period polydisc, where all its rational data are constant.  Throughout
this section, ``very general'' in a connected complex manifold means
outside a countable union of proper closed analytic subsets.

\subsection{Hilbert parameters over an analytic base}

\begin{lemma}[An analytic classifying map]
\label{spec:hilbert-map}
Let $p:\mathcal X\to D$ be a smooth proper holomorphic family of
nonempty, pure $n$-dimensional varieties over a connected complex manifold.
Suppose that an embedding
\[
 \mathcal X\hookrightarrow D\times\mathbb P^N
\]
over $D$ has been given.  There are a reduced complex quasi-projective
scheme $H$, a smooth projective algebraic family
$\mathcal U\to H$, and a holomorphic map
$\eta:D\to H^{\mathrm{an}}$ such that $\mathcal X$ is the embedded
analytic pullback of $\mathcal U$.
\end{lemma}

\begin{proof}
The embedded fibers have a common Hilbert polynomial $P$.
Indeed, Riemann--Roch expresses
$\chi(\mathcal X_t,\mathcal O_{\mathcal X_t}(k))$ in terms of the
fiber integrals of the Chern classes of the relative tangent bundle and
the hyperplane bundle.  These integrals are locally constant in a smooth
proper family, and $D$ is connected.

Let $\operatorname{Hilb}^{P}(\mathbb P^N)$ be the algebraic Hilbert scheme.
Fixed-polynomial boundedness supplies one sufficiently large integer $k$
for which every member $Y$ has
\[
 \dim \operatorname{im}\bigl(H^0(\mathbb P^N,\mathcal O(k))
       \longrightarrow H^0(Y,\mathcal O_Y(k))\bigr)=P(k),
\]
and its degree-$k$ ideal space defines the Hilbert scheme's closed
embedding in the corresponding Grassmannian.  These uniform statements
precede the construction of $\eta$; see
\cite[Theorem~3.2 and Proposition~3.8]{GrothendieckHilbert}.

We check that $t\mapsto H^0(\mathbb P^N,\mathcal I_{\mathcal X_t}(k))$
is holomorphic.  Fix $t_0\in D$.  Since $\mathcal X_{t_0}$ is reduced,
point evaluations separate the restrictions of degree-$k$ polynomials.
Choose finitely many points whose evaluations have rank $P(k)$, and
extend these points to local holomorphic sections of $p$, using its
submersion property.  After locally trivializing $\mathcal O(k)$ along
these sections, evaluation is a holomorphic matrix.  A nonzero
$P(k)$-minor remains nonzero near $t_0$.  Its rank cannot exceed $P(k)$,
because evaluation factors through restriction to $\mathcal X_t$.
Consequently its kernel equals the degree-$k$ ideal space, and these
kernels give the required holomorphic Grassmannian map.

Pointwise this map lies in the Hilbert locus.  Since $D$ is reduced,
the defining equations of that locus vanish on $D$, so the map factors
holomorphically through the Hilbert scheme.  Take $H$ to be the reduction
of the open locus parametrizing smooth pure $n$-dimensional subschemes,
and let $\mathcal U$ be the restricted universal family.  Flatness of the
universal family and smoothness of all its fibers imply that
$\mathcal U\to H$ is smooth.  Its pullback to $D$ and $\mathcal X$
are reduced closed analytic subspaces of $D\times\mathbb P^N$ with the
same fibers, hence the same underlying set.  They therefore coincide.
\end{proof}

For a closed subscheme $Z$ of dimension at most $a$, write
$[Z]_a$ for its $a$-dimensional fundamental cycle, with the lengths at
the generic points as multiplicities.  Components of smaller dimension
do not enter $[Z]_a$.

\begin{lemma}[Classes of flat subschemes]
\label{spec:flat-cycle}
Let $f:M\to R$ be a smooth projective algebraic family of relative
dimension $n$, with $R$ a smooth irreducible complex variety of dimension
$b$.  Let $Z\subset M$ be a closed subscheme, flat over $R$, whose fibers
have dimension $n-2$.  Then
\[
 r\longmapsto
 \operatorname{cl}([Z_r]_{n-2})\in H^4(M_r,\mathbb Q)
\]
is a section of the local system $R^4f^{\mathrm{an}}_*\mathbb Q$.
This includes nonpure and nonreduced fibers.
\end{lemma}

\begin{proof}
The total space $M$ is smooth.  Flatness and the dimension formula give
$\dim Z=n+b-2$, so its top cycle has a class
\[
 \zeta=\operatorname{cl}([Z]_{n+b-2})\in H^4(M,\mathbb Q).
\]
Here the usual cycle class on a smooth variety takes values in ordinary
cohomology; completeness of $M$ is unnecessary.
We claim that, for the inclusion $i_r:M_r\hookrightarrow M$,
\begin{equation}
 \label{spec:restriction}
 i_r^*\zeta=\operatorname{cl}([Z_r]_{n-2}).
\end{equation}

Choose local algebraic coordinates by an \emph{\'etale} map from
a neighborhood $U\subset R$ to $\mathbb A^b$ taking $r$ to $0$.
Shrink $U$ to remove the other points over $0$. The coordinate
functions then cut out a smooth flag
\[
 U=U_0\supset U_1\supset\cdots\supset U_b=\{r\}.
\]
Put $Z_j=Z\times_R U_j$. It remains flat over $U_j$ and has dimension
at most $n+b-2-j$. At each step the next coordinate is a
nonzerodivisor on $U_j$ and hence on the entire scheme $Z_j$.
The Cartier intersection formula therefore gives
\[
 [Z_j]_{n+b-2-j}\longmapsto[Z_{j+1}]_{n+b-3-j},
\]
with the scheme multiplicities
\cite[Lemma~42.29.5, Tag~02TB]{StacksCartier}.
Applying the formula successively yields
\[
 i_r^![Z]_{n+b-2}=[Z_r]_{n-2}.
\]
The formula applies to the whole possibly nonpure or nonreduced
scheme at each cut; no componentwise flatness assertion is needed.
Compatibility of Betti cycle classes with these regular pullbacks proves
\eqref{spec:restriction}.

Finally, the restrictions of a single cohomology class on the total
space of a smooth proper family form a section of its cohomology local
system.  This follows directly from local differentiable triviality.
Apply it to $\zeta$.
\end{proof}

\begin{proposition}[Specialization from very general algebraicity]
\label{spec:specialization}
Let $D$ be a connected polydisc, and let
$p:\mathcal X\to D$ be a smooth proper holomorphic family of
$n$-dimensional varieties, $n\geq2$, with a relative embedding in
$D\times\mathbb P^N$.  Let
\[
 \theta\in H^0(D,R^4p_*\mathbb Q)
\]
be a flat rational class.  If $\theta_t$ is the class of a rational
codimension-two algebraic cycle for very general $t\in D$, then the same
holds for every $t\in D$.
\end{proposition}

\begin{proof}
Fix a desired point $0\in D$ and trivialize $R^4p_*\mathbb Q$ over $D$.
Use Lemma~\ref{spec:hilbert-map} to write the family as the pullback of
$\mathcal U\to H$ along $\eta:D\to H^{\mathrm{an}}$.
For each possible Hilbert polynomial $Q$ of degree $n-2$, consider
\[
 T_Q=\operatorname{Hilb}^{Q}(\mathcal U/H).
\]
This is projective over $H$ and has a flat universal subscheme.
It has finitely many reduced irreducible components.  Resolve each
component by a projective birational morphism, obtaining a finite
collection of smooth irreducible varieties $R\to (T_Q)_{\mathrm{red}}$
whose images cover $(T_Q)_{\mathrm{red}}$; use
\cite{Hironaka}.  Thus each composite $R\to H$ is proper.
Pullback gives a smooth projective ambient family $\mathcal U_R\to R$
and a flat closed subscheme $Z_R\subset\mathcal U_R$.  By
Lemma~\ref{spec:flat-cycle}, the classes of the top cycles of its fibers
are sections of $R^4(\mathcal U_R/R)_*\mathbb Q$.

Form the analytic spaces
\[
 K_R=D\times_{H^{\mathrm{an}}}R^{\mathrm{an}}.
\]
They are proper over $D$.  Their pulled-back ambient families agree with
the pullback of $\mathcal X$, so the preceding cycle classes become
locally constant functions on $K_R$ with values in the fixed marked
space $H^4(\mathcal X_0,\mathbb Q)$.  The spaces $K_R$ may be singular
or nonreduced: this causes no difficulty, since local constancy was
established on the smooth parameter $R$ before base change.

There are countably many $Q$ and finitely many choices of $R$ for each
$Q$.  Each $K_R$ is second countable and has a locally finite family of
closed irreducible analytic components.  It consequently has at most
countably many such components.  Every component is proper over $D$;
Remmert's Proper Mapping Theorem makes its image a closed analytic
subset of $D$ \cite[\S7, Satz~1]{Grauert}.  Let $E$ be the union of the images
of all components whose image is not all of $D$.  Thus $E$ is a
countable union of proper closed analytic subsets.

Choose $t\in D$ outside $E$ and outside the exceptional sets in the
hypothesis.  Such a point exists by the Baire Category Theorem.  Choose
an actual cycle representing $\theta_t$, and write it as a finite sum
\[
 \gamma_t=\sum_{j=1}^{s}a_j[Y_{j,t}],
 \qquad a_j\in\mathbb Q,
\]
where each $Y_{j,t}\subset\mathcal X_t$ is an integral subvariety of
codimension two.  Its Hilbert point lifts to one of the parameters $R$,
and hence to a point $u_j\in K_R$ above $t$.  Choose an irreducible
component $C_j$ containing $u_j$.  Because $t\notin E$, the map
$C_j\to D$ is surjective.  In particular, there is a point
$v_j\in C_j$ above $0$.

An irreducible complex analytic space is connected.  The marked cycle
class is therefore constant on $C_j$, and the fiber of the universal
subscheme at $v_j$ has top cycle $\delta_{j,0}$ satisfying
\[
 \operatorname{cl}(\delta_{j,0})
 =\operatorname{cl}([Y_{j,t}])
 \quad\text{under the marking.}
\]
These are algebraic cycles on $\mathcal X_0$: the fibers of the
universal subscheme are algebraic in $\mathbb P^N$, and the embedded
fiber identification is algebraic by GAGA \cite{SerreGAGA}.
Consequently
\[
 \gamma_0=\sum_{j=1}^{s}a_j\delta_{j,0}
 \quad\text{satisfies}\quad
 \operatorname{cl}(\gamma_0)=\theta_0.
\]
The argument applies to any chosen point of $D$.
\end{proof}

The avoidance in this proof is essential: density of fibers with
algebraic classes by itself would not ensure that the components
parametrizing those cycles reach the desired fiber.  Proposition~\ref{spec:specialization}
requires actual algebraicity at very general points.  It assumes no
principle asserting algebraicity from the Hodge property of a flat
class.

\subsection{A marked family for the primitive Clifford tensor}

\begin{lemma}[Local polarized families]
\label{spec:local}
Fix a primitively amply polarized K3 surface $(S_0,h)$.
After choosing a sufficiently small full-dimensional polarized period
polydisc $D$ about its period, there are smooth proper holomorphic
families
\[
 \mathcal S\longrightarrow D,
 \qquad \mathcal A\longrightarrow D,
\]
with the following properties.
The K3 family has central fiber $(S_0,h)$, a relative ample line bundle,
and a marking of its integral second cohomology.
The primitive rational local system is a fixed quadratic space
$(V,q)$ of dimension $21$.
The abelian family has marked rational first cohomology
$W=C^+(V,q)$ with the full Clifford Hodge structures.
Fixed rational polarization and parity-identification data give a flat
family of standard tensor maps
\[
 \mu_t:V\longrightarrow H^1(\mathcal A_t,\mathbb Q)
                       \otimes H^1(\mathcal A_t,\mathbb Q).
\]
Finally, after shrinking $D$, the product
\[
 \mathcal X=\mathcal S\times_D\mathcal A\times_D\mathcal A
\]
admits a relative embedding in $D\times\mathbb P^N$.
\end{lemma}

\begin{proof}
Polarized deformation theory and local Torelli give the smooth K3
family over a $19$-dimensional polydisc on which the period map is a
local isomorphism; see \cite[Chapter~6, Section~2]{Huybrechts}.
Since $D$ is contractible, its cohomology local systems are marked.
The class $h$ remains of type $(1,1)$.  We recall why it is represented
by a line bundle on the total family.
Local differentiable triviality over the contractible base extends
$h$ to $H^2(\mathcal S,\mathbb Z)$.  The polydisc is Stein, and coherent
Leray and Cartan's Theorem~B identify
\[
 H^2(\mathcal S,\mathcal O_{\mathcal S})
 =H^0(D,R^2p_*\mathcal O_{\mathcal S}),
 \qquad p:\mathcal S\to D.
\]
The direct image on the right is locally free and commutes with base
change.  The image of $h$ in it vanishes fiberwise, hence vanishes.
The exponential sequence now gives a holomorphic line bundle
$\mathcal L$ with first Chern class $h$.

A high power of $\mathcal L$ is very ample on $S_0$ and has vanishing
higher cohomology there.  Proper cohomology and base change extend a
basis of its sections over a neighborhood of $0$.  They give a
relative embedding after shrinking: basepoint freeness and immersion
persist near the central fiber, and properness excludes pairs of
distinct points with the same image.  For the latter assertion, any
sequence of such pairs approaching the central fiber either limits
to distinct points, contrary to its embedding, or to a single point,
contrary to the local embedding furnished by the immersion.
This proves local relative projectivity, and also that $h$ stays ample.
The analytic tools used here are the proper direct-image and base-change
theorems and the Stein vanishing theorem; see
\cite[\S6, Hauptsatz~I; \S7, S\"atze~3,5; \S2, Satz~3]{Grauert}.

On the fixed Clifford space $W$, the Hodge filtration and rational
alternating polarization are those constructed in
Proposition~\ref{cliff:construction}.  To recall the holomorphic dependence,
if $\omega_t$ spans the period line, then
\[
 F^1W_{\mathbb C}
 =\operatorname{im}\bigl(L_{\omega_t}:C^-(V_{\mathbb C},q)
                                      \longrightarrow C^+(V_{\mathbb C},q)\bigr).
\]
This has constant half dimension and varies holomorphically.
The polarization can be fixed on a sufficiently small chart.
Choose a rational symplectic basis and let its integral span be a
lattice $\Lambda$ in $W$. It is unimodular for the chosen symplectic form and
does not change the rational Hodge structure.
Since $W$ is cohomology, the associated torus is explicitly
\[
 \mathcal A_t=(F^1W_t)^\vee/\operatorname{image}(\Lambda^\vee),
\]
where a lattice functional is restricted to $F^1W_t$.
Thus $H_1(\mathcal A_t,\mathbb Z)=\Lambda^\vee$,
$H^1(\mathcal A_t,\mathbb Z)=\Lambda$, and
$H^{1,0}(\mathcal A_t)=F^1W_t$, compatibly with the markings.
The polarization supplies the Riemann form on the dual lattice.
Equivalently, the marked filtrations give a holomorphic map to Siegel
space and $\mathcal A\to D$ is the pullback of its marked torus family.
The full markings supply, in particular, level-$3$ data. The spin
action is defined on the fixed Clifford space over the polydisc,
so there is no monodromy lifting problem on this base.

For completeness, in symplectic period coordinates this torus family
is
\[
 \mathbb C^g/(\mathbb Z^g+\tau(t)\mathbb Z^g),
 \qquad \tau(t)^{\mathsf t}=\tau(t),\quad
          \operatorname{Im}\tau(t)>0.
\]
The factors
\[
 \exp\bigl(-\pi i\,n^{\mathsf t}\tau(t)n
               -2\pi i\,n^{\mathsf t}z\bigr),
 \qquad m+\tau(t)n\in\mathbb Z^g+\tau(t)\mathbb Z^g,
\]
define a holomorphic theta line bundle, ample on every fiber.
The preceding extension-of-sections argument gives a relative
projective embedding after shrinking $D$.

Fix a rational nonisotropic vector $z_0\in V$ and the
polarization identification $j:W^\vee(-1)\xrightarrow{\sim}W$.
The maps
\[
 V(1)\longrightarrow\operatorname{End}(W),
 \qquad x\longmapsto\bigl(w\longmapsto xwz_0\bigr)
\]
give $\mu_t$ after untwisting and applying $j$ to the dual slot.
Clifford multiplication, $z_0$, and $j$ are constant rational data in
the markings.  Thus $\mu_t$ is flat as a rational tensor and is a Hodge
map in every fiber.  A Segre embedding combines the relative embeddings
of $\mathcal S$ and the two copies of $\mathcal A$.
\end{proof}

\begin{corollary}[The primitive map at every polarized fiber]
\label{spec:primitive}
Fix a polarized moduli component, and assume that at very general
points the nonzero algebraic input of Theorem~\ref{main:generic} holds.
Then at every $(S_0,h)$ in that component the standard full primitive
Clifford tensor is induced by a rational codimension-two algebraic
cycle.  The same conclusion holds if the hypothesis is made outside
countably many proper closed analytic subsets in a full-dimensional
polarized period chart about each desired point.
\end{corollary}

\begin{proof}
Take the families in Lemma~\ref{spec:local} about $(S_0,h)$.
Under the markings, extend $\mu_t$ by zero on $\mathbb Qh$ and on the
other degrees of $H^*(\mathcal S_t,\mathbb Q)$.
Poincar\'e duality and the K\"unneth decomposition give a unique flat
class
\[
 \theta\in H^0(D,R^4(\mathcal X/D)_*\mathbb Q)
\]
with this action.  More explicitly, for a rational basis $(v_i)$ of
$V$, let $(q^{ij})$ be the inverse cup-pairing matrix.  Then
\begin{equation}
 \label{spec:theta}
 \theta_t=\sum_{i,j}q^{ij}v_i\otimes\mu_t(v_j)
 \ \in\ V\otimes
 \bigl(H^1(\mathcal A_t)\otimes H^1(\mathcal A_t)\bigr)
 \ \subset\ H^4(\mathcal X_t,\mathbb Q).
\end{equation}

By Proposition~\ref{gen:genericity}, after excluding countably many
proper closed analytic subsets we have
$\operatorname{Hg}(V_t)=\operatorname{SO}(V,q)$.
At a point satisfying this condition and the assumed input,
Theorem~\ref{gen:generation} produces a cycle inducing $\mu_t$ on $V_t$.
Its full class can be made exactly $\theta_t$.
Indeed, on $\mathcal S_t$ the algebraic correspondence
\begin{equation}
 \label{spec:primitive-projector}
 \pi_{V,t}=[\Delta_{\mathcal S_t}]
 -[o_t]\times\mathcal S_t-\mathcal S_t\times[o_t]
 -\frac{h\times h}{\int_{\mathcal S_t}h^2},
 \qquad o_t\in\mathcal S_t,
\end{equation}
acts identically on $V_t$ and kills $H^0$, $H^4$, and $\mathbb Qh$.
There is no odd K3 cohomology.  Precompose the chosen cycle with
$\pi_{V,t}$ and postcompose with the algebraic projector onto
$H^1(\mathcal A_t)\otimes H^1(\mathcal A_t)$, as in
Lemma~\ref{cyc:projectors}.  These operations preserve codimension two.
The resulting cycle has the prescribed action in every source degree,
so Poincar\'e duality and K\"unneth identify its entire class with
\eqref{spec:theta}.  The point $o_t$ is needed only on this fiber; its
cohomology class, and hence the projector's cohomological action, is
canonical.

If the hypothesis is stated in algebraic moduli, its exceptional loci
pull back to proper closed analytic subsets of $D$.  They remain proper
because the full-dimensional polarized period chart has open image
locally in that component, up to a finite automorphism quotient.
Thus $\theta_t$ is algebraic very generally on $D$ in either version
of the hypothesis.  Proposition~\ref{spec:specialization} produces an
actual rational codimension-two cycle with class $\theta_0$.
Its action on $V$ is exactly $\mu_0$. Undoing the fixed invertible right multiplier on the output Hom slot, using Lemma~\ref{cliff:slots}, also algebraically realizes $m^q$ itself.
\end{proof}

Neither the original witness abelian varieties nor their cycles have
to form a family.  Theorem~\ref{gen:generation} first transfers the
input to the fixed local full Clifford family.  The specialization
argument then uses a finite representative of one sufficiently general
cycle, keeping its rational coefficients unchanged.

\section{Restriction and the prescribed data at a special fiber}
\label{data:section}

Let $(S,h)$ now be any fiber under consideration, and put
\[
 V=h^\perp\subset H^2(S,\mathbb Q),\qquad
 T=\operatorname{NS}(S)_{\mathbb Q}^{\perp}\subset V.
\]
Suppose that the primitive Clifford multiplication for the cup form
$q$ on $V$ has been algebraically realized. We show that this gives
the stipulated standard map on $T$, including its chosen rational
quadratic polarization. The arguments use no hypothesis on the
Hodge group at this fiber.

\begin{proposition}[Restriction to the transcendental part]
\label{data:restriction}
Set $q_0=q|_T$. Algebraic realization of $m^q$ on $V$ implies
algebraic realization of $m^{q_0}$ on $T$.
\end{proposition}

\begin{proof}
The orthogonal decomposition $V=T\oplus T^\perp$ is a decomposition
of rational Hodge structures, and $T$ contains the period plane.
By the Clifford Basis Theorem there are inclusions
\[
 C^\pm(T,q_0)\hookrightarrow C^\pm(V,q).
\]
They are weight-one Hodge inclusions: the element $J$ defining the
left Hodge circle already belongs to $C^+(T_{\mathbb R},q_0)$.
Let $i:C^+(T,q_0)\to C^+(V,q)$ be the inclusion and choose a rational
Hodge retraction $p:C^-(V,q)\to C^-(T,q_0)$. Such a retraction
exists by polarizability: the orthogonal complement for a
weight-one polarization is a rational Hodge complement.
For every $x\in T$,
\begin{equation}\label{data:restriction-identity}
 p\,L_x^{V}\,i=L_x^{T},
\end{equation}
because $xC^+(T,q_0)\subset C^-(T,q_0)$ and $p$ restricts to the
identity there. Lemma~\ref{cliff:slots} realizes the two operations
on the abelian factors. Restricting the resulting cycle action to
$T$ gives the asserted equality; no operation on the K3 factor
is needed.
\end{proof}

\begin{proposition}[Change of quadratic polarization]
\label{data:polarization}
Let $q_0$ and $q_1$ be rational quadratic Hodge polarizations on
the same K3-type Hodge structure $T$, allowing either overall sign
as in Section~\ref{cliff:section}. If $T\subset H^2(S,\mathbb Q)$
and $m^{q_0}$ is algebraically realized, then so is $m^{q_1}$.
\end{proposition}

\begin{proof}
\emph{Step 1: square-root isometries.}
Write
\begin{equation}\label{data:P}
 q_1(u,v)=q_0(Pu,v).
\end{equation}
The two form isomorphisms from $T$ to its twisted dual show that
$P$ is an invertible rational Hodge endomorphism. Symmetry of the
forms makes $P$ self-adjoint for $q_0$.
Let $K$ be $+1$ on the real period plane and $-1$ on the real
$(1,1)$-part, and let $\epsilon_0$ be the sign of $q_0$ on the
period plane. Then
\[
 (u,v)_0=\epsilon_0 q_0(Ku,v)
\]
is positive definite. Since $P$ commutes with $K$, it is also
self-adjoint for this positive definite real form. Hence $P$ is
semisimple with real nonzero eigenvalues. Its distinct eigenvalues
$\lambda_1,\ldots,\lambda_r$ and spectral projectors
$p_1,\ldots,p_r$ are defined over $\overline{\mathbb Q}$, the latter
being polynomials in $P$.

Choose $\nu_i\in\overline{\mathbb Q}$ with $\nu_i^2=\lambda_i$.
For each $s=(s_1,\ldots,s_r)\in\{1,-1\}^r$, put
\begin{equation}\label{data:roots}
 R_s=\sum_{i=1}^r s_i\nu_i p_i.
\end{equation}
These are invertible self-adjoint operators with $R_s^2=P$.
Consequently
\[
 q_0(R_su,R_sv)=q_0(Pu,v)=q_1(u,v).
\]
Thus $R_s:(T,q_1)_{\overline{\mathbb Q}}\to
(T,q_0)_{\overline{\mathbb Q}}$ is an isometry and extends to a
parity-preserving Clifford algebra isomorphism
\[
 F_s:C(T,q_1)_{\overline{\mathbb Q}}
          \xrightarrow{\sim}C(T,q_0)_{\overline{\mathbb Q}}.
\]
\smallskip\noindent
\emph{Step 2: weight-one Hodge maps.}
We show that the parity restrictions of $F_s$ and their inverses lie
in the $\overline{\mathbb Q}$-span of rational weight-one Hodge maps.

Let $\mathcal J$ be the joint rational Hodge group of $T$ and the
four weight-one spaces $P_j^\pm=C^\pm(T,q_j)$, $j=0,1$.
Here the joint Hodge group is the smallest rational algebraic
subgroup containing their simultaneous real Hodge circle.
Its geometric identity component is defined over $\mathbb Q$ and
already contains that connected circle. Minimality therefore gives
$\mathcal J=\mathcal J^\circ$, so $\mathcal J$ is geometrically
connected. Write $\rho_T$ for its action on
$T$. Since $P$ is a rational Hodge endomorphism, it commutes with
$\rho_T(\mathcal J)$; therefore every $R_s$ does as well after
scalar extension.

For each $j$, the projection of $\mathcal J$ to the actions on
$(T,P_j^+,P_j^-)$ lies in the image of
\[
 \operatorname{Spin}(T,q_j)\longrightarrow
 \operatorname{GL}(T)\times\operatorname{GL}(P_j^+)
                      \times\operatorname{GL}(P_j^-),
\]
using the orthogonal action on $T$ and left multiplication on
the parity spaces. Indeed, this rational algebraic image contains
the projected circle by \eqref{cliff:circle}. The representation
is faithful, since the left action on $C^+$ recovers a spin
element by evaluation at $1$. Thus it identifies the spin group
with its algebraic image, and gives algebraic homomorphisms
\[
 \rho_j:\mathcal J\longrightarrow\operatorname{Spin}(T,q_j)
\]
covering the same action $\rho_T$.

The isometry $R_s$ induces an isomorphism of spin groups
$\widehat R_s:\operatorname{Spin}(T,q_1)_{\overline{\mathbb Q}}
\to\operatorname{Spin}(T,q_0)_{\overline{\mathbb Q}}$.
Since $R_s$ commutes with $\rho_T$, both
$\widehat R_s\rho_1$ and $\rho_0$ lift the same orthogonal
representation of $\mathcal J_{\overline{\mathbb Q}}$.
Their pointwise quotient therefore takes values in the central
kernel $\mu_2$ of the spin covering. This quotient is a morphism
from a connected group to a finite group, and is the identity at
the identity. Hence
\begin{equation}\label{data:lift-agreement}
 \widehat R_s\rho_1=\rho_0.
\end{equation}
Because $F_s$ is an algebra isomorphism,
\eqref{data:lift-agreement} says precisely that $F_s$ intertwines
the left $\mathcal J$-actions on both parities.

For finite-dimensional rational representations, equivariant
Hom-spaces commute with field extension: equivariance is a set
of rational linear equations on the matrix of a map. Rational
maps intertwine $\mathcal J$ exactly when they intertwine its
Hodge circle. Thus
\begin{equation}\label{data:hom-descent}
 \operatorname{Hom}_{\mathcal J}(P_1^\pm,P_0^\pm)
       \otimes_{\mathbb Q}\overline{\mathbb Q}
 =\operatorname{Hom}_{\mathcal J_{\overline{\mathbb Q}}}
       (P_{1,\overline{\mathbb Q}}^\pm,
        P_{0,\overline{\mathbb Q}}^\pm)
\end{equation}
expresses $F_s^\pm$, and likewise $(F_s^\pm)^{-1}$, as finite
linear combinations of rational weight-one Hodge maps.

\smallskip\noindent
\emph{Step 3: recovery and rational descent.}
Fix abelian realizations and rational dualities for the $q_1$
target, and let $\mathcal A$ be the rational subspace of maps
\[
 T\longrightarrow P_1^-\otimes P_1^+
\]
obtained by restricting rational codimension-two cycle actions
from $S$ and projecting to this K\"unneth component. It is a
subspace of a finite-dimensional rational Hom-space.
By Lemma~\ref{cliff:slots} and \eqref{data:hom-descent}, conjugating
the algebraically realized $q_0$ multiplication on its Hom slots
gives elements of $\mathcal A\otimes\overline{\mathbb Q}$.
The exact Clifford identity is
\begin{equation}\label{data:conjugation}
 (F_s^-)^{-1}L_x^{q_0}F_s^+
       =L_{R_s^{-1}x}^{q_1}\qquad(x\in T).
\end{equation}
Here the conjugation is expanded using rational Hodge maps on
the two abelian factors; it does not apply $R_s^{-1}$ to the
cohomology of $S$ by a correspondence.

An elementary sign average recovers the unmodified source:
\begin{equation}\label{data:sign-average}
 p_i=\frac{\nu_i}{2^r}
        \sum_{s\in\{1,-1\}^r}s_iR_s^{-1},
 \qquad \operatorname{id}_T=\sum_{i=1}^r p_i.
\end{equation}
Indeed $R_s^{-1}=\sum_j s_j\nu_j^{-1}p_j$ and
$\sum_s s_i s_j=2^r\delta_{ij}$.
Equations~\eqref{data:conjugation} and \eqref{data:sign-average}
therefore place the tensor realization of $x\mapsto L_x^{q_1}$
in $\mathcal A\otimes\overline{\mathbb Q}$.
This map is rational, so Lemma~\ref{cyc:descent} places it in
$\mathcal A$ itself. A finite rational linear combination of
actual cycles realizes it.
\end{proof}

\begin{remark}\label{data:scope}
The proof never uses an algebraic correspondence representing
$P$, any $p_i$, or any $R_s$ on the K3 surface. They enter only
the identities that produce weight-one intertwiners and the
finite rational descent. In particular $P$ need not be scalar.
The argument applies with additional real multiplication or
complex multiplication endomorphisms, and requires no description of the special
Hodge group. It applies also in ranks $2$, $3$, and $4$.
In rank $2$, the full even-Clifford space is two-dimensional and
$\operatorname{Spin}_2\to\operatorname{SO}_2$, over an algebraic
closure, is the squaring map on $\mathbb G_m$ with kernel
$\mu_2$. The faithful left Clifford action and the connected-lift
argument remain valid. Higher small ranks introduce no exception.
We make no assertion that all Hodge embeddings for a smaller
special Hodge group are generated by these operations.

The opposite-sign case is directly visible in rank two. Take
$q_0=\operatorname{diag}(1,1)$ and $q_1=-q_0$ in the same period
basis. Then $J_1=-e_1e_2$ and the isometry $R=i\operatorname{id}$
induces $F(J_1)=-(ie_1)(ie_2)=J_0$. Thus the nonreal square root
preserves the specified weight-one circle, as the lift argument predicts.
\end{remark}

For completeness, a rational Hodge automorphism $D$ of $T$ in a
prescribed identification can also be incorporated without
applying it as a K3 correspondence. First use
Proposition~\ref{data:polarization} for the pullback form $D^*q_1$.
The rational isometry $D:(T,D^*q_1)\to(T,q_1)$ induces
weight-one Hodge isomorphisms on Clifford spaces: the same
connected-lift argument proves this, since $D$ commutes with the
joint Hodge action. Transport on the two Hom slots then gives
exactly $x\mapsto L_{Dx}^{q_1}$.

Finally apply Proposition~\ref{data:restriction}, then
Proposition~\ref{data:polarization}, and then the prescribed
right multiplier and duality of \eqref{cliff:standard}.
Lemma~\ref{cliff:slots} transports the resulting tensor to the
actual $H^1(A_S,\mathbb Q)$ model. The target K\"unneth projectors
give a cycle
\[
 \Gamma_S\in\operatorname{CH}^2(S\times A_S\times A_S)_{\mathbb Q}
 \quad\text{with}\quad
 (\Gamma_S)_*x=\kappa_S(x)\quad\text{for every }x\in T.
\]
This is equality with the prescribed full map, including its
rational normalization, and completes the exact-data step in
Theorems~\ref{main:generic} and~\ref{main:specialization}.

\section{Completion of the reductions}\label{sec:conclusion}

\begin{proof}[Proof of Theorem~\ref{main:generic}]
Proposition~\ref{gen:genericity} gives $T=V$.
Proposition~\ref{cyc:equivalence} proves the equivalence of
conditions \textup{(i)} and \textup{(ii)}.
From condition \textup{(i)}, Theorem~\ref{gen:generation}
realizes the full primitive Clifford multiplication.
Proposition~\ref{data:polarization} and
Lemma~\ref{cliff:slots} then give condition \textup{(iii)}
with the stipulated data.
Conversely, condition \textup{(iii)} gives condition \textup{(i)}
by taking $B=A_S\times A_S$ and $\alpha=\Gamma_S$.
\end{proof}

\begin{proof}[Proof of Theorem~\ref{main:specialization}]
By Proposition~\ref{cyc:equivalence}, either assumed input gives
the nonzero correspondence very generally.
Corollary~\ref{spec:primitive}, followed by the inverse right
multiplier in Lemma~\ref{cliff:slots}, algebraically realizes
$m^q$ on $V$ at the desired point.
Proposition~\ref{data:restriction} gives $m^{q|_T}$ on $T$.
Proposition~\ref{data:polarization} and
Lemma~\ref{cliff:slots} then give the prescribed quadratic form,
parity identification, duality, and isogenous realization,
hence \eqref{eq:cycle-target}.
The operations on the abelian factors preserve the degree shift,
so the resulting cycle has codimension two.
\end{proof}

The conclusions include Picard jumps, extra real or complex multiplication, and the smallest possible transcendental ranks. They require the prescribed multiplication tensor at those points; they do not assert generation of every Hodge embedding for a smaller special Hodge group. The mathematical content of the reduction is conditional: producing a very-general nonzero algebraic correspondence, or equivalently the surface cover of Theorem~\ref{main:generic}, remains the external existence problem.

\appendix
\section{Established inputs}\label{app:dependencies}
The following background results are the external dependencies.  The
representation-width argument, recovery by independent tensor operations,
change of quadratic form, and specialization of the prescribed class are
proved in the text.  No algebraicity theorem for the Kuga--Satake class is
an input.  Page numbers for \cite{Huybrechts} refer to the linked draft.

\begingroup
\small
\setlength{\tabcolsep}{5pt}
\renewcommand{\arraystretch}{1.13}
\begin{longtable}{@{}>{\raggedright\arraybackslash}p{.40\textwidth}>{\raggedright\arraybackslash}p{.31\textwidth}>{\raggedright\arraybackslash}p{.23\textwidth}@{}}
\toprule
Established statement and hypotheses & Reference and locator & Role \\
\midrule
\endfirsthead
\toprule
Established statement and hypotheses & Reference and locator & Role \\
\midrule
\endhead
\bottomrule
\endfoot

For smooth projective complex varieties, Hodge decomposition, Poincar\'e
duality, the Lefschetz theorem on $(1,1)$-classes, hard Lefschetz, and the
Hodge--Riemann relations hold; pullback, Gysin maps, and products have their
usual Hodge types. &
\cite{Voisin}, \S\S6.1--6.3, pp.~139--154;
\S\S7.1--7.3, pp.~157--182. &
Pairings, primitive cohomology, divisor classes, and nonvanishing
calculations. \\[3pt]

Polarizable rational Hodge structures have Hodge complements, and their
rational Hodge group is reductive.  The same
holds for a finite joint collection. &
\cite{Deligne}, Proposition 3.6 and its proof, electronic p.~27;
\cite{Voisin}, \S7.3.1. &
Rational splittings and the joint Hodge-group Lie algebra. \\[3pt]

In characteristic zero algebraic group homomorphisms have closed image
and surjective differential onto its Lie algebra.  A unipotent group has
nonzero invariants in every nonzero representation.  A reductive Lie
algebra is the sum of its center and simple ideals. &
\cite{Borel}, Corollary 1.4, Theorem 4.8, \S7,
and Proposition 14.2;
\cite{FultonHarris}, Lectures 9 and 14. &
The rational simple ideal above the orthogonal quotient and the
commuting kernel. \\[3pt]

A complex K3 surface has the rank-$22$ cup lattice and Hodge numbers used
here.  Its universal deformation is smooth of dimension $20$; the local
period map is an isomorphism.  Fixing a nonzero integral $(1,1)$-class
cuts out the corresponding local period hypersurface and extends its
line bundle. &
\cite{Huybrechts}, Chapter 1, \S\S2--3;
Chapter 6, Corollary 2.7 and Proposition 2.8, p.~107;
\S2.4, p.~108. &
Marked polarized deformation charts and the rank-$21$ primitive space. \\[3pt]

The full even Clifford construction of a polarized Hodge structure of K3
type is polarizable of weight one.  Clifford multiplication with an
invertible right multiplier gives the standard Hodge embedding. &
\cite{Huybrechts}, Chapter 4, Propositions 2.5--2.6 and Remark 2.7,
pp.~66--68;
\cite{vanGeemen}, Propositions 5.9 and 6.3. &
The prescribed full target, Hodge types, and fixed construction data. \\[3pt]

Complex semisimple Lie algebras have completely reducible finite-dimensional
representations and highest-weight classification.  In type $B_n$ the
dominant-weight coordinates and spin weights have the stated forms.
The Clifford Basis Theorem and spin double covering hold in characteristic
zero. &
\cite{FultonHarris}, Lectures 14 and 18--20;
\S19.4, pp.~294--296;
\S\S20.1--20.2, pp.~299--312. &
The width bound, spin multiplicity calculation, and finite-kernel lift
comparison. \\[3pt]

The functor $H^1$ identifies complex abelian varieties up to isogeny
contravariantly with polarizable rational weight-one Hodge structures.
A polarized complex torus is projective. &
\cite{Milne}, Theorems 6.7--6.8 and Corollary 6.9,
pp.~72--73 (stated covariantly for homology). &
Algebraic realization of every rational weight-one slot map and of the
Clifford spaces. \\[3pt]

For a smooth projective complex variety $Y$, its Albanese is an abelian
variety, and its Albanese map induces an isomorphism on holomorphic
one-forms. &
\cite{Albanese}, Chapter VI, \S9.2, pp.~315--316,
formulas (9.11)--(9.13). &
Conversion of the surface-cover input into a correspondence involving an
abelian variety. \\[3pt]

On smooth projective varieties, correspondences compose by
pullback, intersection, and proper pushforward, with the indicated
codimension.  Cycle classes commute with these operations. &
\cite{Fulton}, \S16.1, pp.~305--315, and
\S\S19.1--19.2, pp.~371--385;
\cite{StacksCorrespondences}, Lemmas 45.3.1 and 45.3.6. &
Actual rational cycles throughout, with controlled degrees and graph
transposes. \\[3pt]

For a projective family over a noetherian base, each fixed-polynomial
relative Hilbert scheme is projective and has a flat universal
subscheme.  One uniform sufficiently high degree gives its Grassmannian
realization. &
\cite{GrothendieckHilbert}, Theorem 3.2, p.~260;
Proposition 3.8, p.~265. &
Countably many proper cycle parameter spaces and the analytic
classifying map. \\[3pt]

Flat morphisms of finite presentation with smooth fibers are smooth.
In a proper flat family the smooth-fiber locus is open: remove the
proper image of the closed nonsmooth locus.  High powers of an ample
line bundle on a projective variety are very ample and have vanishing
higher cohomology.  Riemann--Roch expresses their Euler characteristics
by Chern classes. &
\cite{StacksFamilies}, tags 01V4, 01V8, 01VU, 0B5U;
\cite{Fulton}, \S18.3. &
The smooth ambient Hilbert family, projective embeddings, and the fixed
Hilbert polynomial. \\[3pt]

If a Cartier equation is a nonzerodivisor on a closed subscheme $Z$ of
dimension at most $k+1$, its Cartier pullback takes $[Z]_{k+1}$ to the
$k$-dimensional cycle of the cut, with scheme multiplicities. &
\cite{StacksCartier}, Lemma 42.29.5(1), tag 02TB;
\cite{Fulton}, Chapters 2 and 6. &
Specialization of top cycles, including nonreduced and nonpure Hilbert
fibers. \\[3pt]

A proper holomorphic map has coherent higher direct images.  For a
coherent sheaf flat over a reduced base, fiber cohomology is upper
semicontinuous; constant dimensions give local freeness and fiber base
change.  Coherent sheaves on Stein spaces have vanishing higher
cohomology. &
\cite{Grauert}, \S6, Hauptsatz I, p.~59;
\S7, S\"atze 3 and 5, pp.~62, 64;
\S2, Satz 3, p.~17, with the 1963 correction. &
Extension of polarizations, local relative projective embeddings, and
holomorphic Grassmannian maps. \\[3pt]

A proper smooth submersion is locally differentiably trivial.  Proper
holomorphic images of closed analytic sets are analytic.  Irreducible
components of a complex analytic set are closed and locally finite;
in a second-countable space there are only countably many. &
\cite{Ehresmann}, Lemma 9.2, p.~77;
\cite{Grauert}, \S7, Satz 1, p.~60;
\cite{Albanese}, Chapter II, Theorem 5.3, pp.~102--103. &
Marked flat cohomology, proper-image avoidance, and transport to the
chosen fiber. \\[3pt]

Complex quasi-projective varieties admit projective resolutions.
Closed analytic subvarieties of projective space and holomorphic maps
between projective varieties are algebraic.  Projective coherent
cohomology agrees with analytic cohomology. &
\cite{Hironaka}, Part I, Main Theorems I and I$^*$;
\cite{SerreGAGA}, Theorems 1--3, pp.~19--20,
and Propositions 13--15, pp.~29--30. &
Smooth cycle parameters and surfaces; algebraicity of the resulting
fiber cycles and weight-one homomorphisms. \\

\end{longtable}
\endgroup

\begin{thebibliography}{99}

\bibitem{Borel}
A. Borel,
\emph{Linear Algebraic Groups}, second edition,
Graduate Texts in Mathematics 126, Springer, 1991.
Corollary 1.4, Theorem 4.8, \S7, and Proposition 14.2.
\url{https://doi.org/10.1007/978-1-4612-0941-6}.

\bibitem{Deligne}
P. Deligne,
Hodge cycles on abelian varieties (notes by J.~S. Milne),
in \emph{Hodge Cycles, Motives, and Shimura Varieties},
Lecture Notes in Mathematics 900, Springer, 1982, pp.~9--100.
Proposition 3.6; revised electronic text (2003), p.~27:
\url{https://www.jmilne.org/math/Documents/Deligne82_2003.pdf}.

\bibitem{Albanese}
J.-P. Demailly,
\emph{Complex Analytic and Differential Geometry},
online book, version of 21 June 2012.
Chapter II, Theorems 5.3 and 8.8; Chapter VI, \S9.2.
\url{https://people.math.harvard.edu/~demarco/Math274/Demailly_ComplexAnalyticDiffGeom.pdf}.

\bibitem{Fulton}
W. Fulton,
\emph{Intersection Theory}, second edition,
Ergebnisse der Mathematik und ihrer Grenzgebiete (3) 2,
Springer, 1998.
Chapters 2 and 6; \S\S16.1, 18.3, 19.1--19.2.
\url{https://doi.org/10.1007/978-1-4612-1700-8}.

\bibitem{FultonHarris}
W. Fulton and J. Harris,
\emph{Representation Theory: A First Course},
Graduate Texts in Mathematics 129, Springer, 1991.
Lectures 14, 18--20; in particular \S19.4 and \S\S20.1--20.2.
\url{https://doi.org/10.1007/978-1-4612-0979-9}.

\bibitem{vanGeemen}
B. van Geemen,
Kuga--Satake varieties and the Hodge conjecture,
in \emph{The Arithmetic and Geometry of Algebraic Cycles}
(Banff, 1998), NATO Science Series C 548,
Kluwer, 2000, pp.~51--82.
Propositions 5.9 and 6.3; preprint version:
\url{https://arxiv.org/abs/math/9903146}.

\bibitem{Grauert}
H. Grauert,
Ein Theorem der analytischen Garbentheorie und die Modulr\"aume
komplexer Strukturen,
\emph{Publ. Math. Inst. Hautes \'Etudes Sci.} \textbf{5} (1960), 5--64;
correction, \emph{ibid.} \textbf{16} (1963), 35--36.
\url{https://www.numdam.org/item/PMIHES_1960__5__5_0/};
\url{https://www.numdam.org/item/PMIHES_1963__16__35_0/}.

\bibitem{GrothendieckHilbert}
A. Grothendieck,
Techniques de construction et th\'eor\`emes d'existence en g\'eom\'etrie
alg\'ebrique IV: les sch\'emas de Hilbert,
\emph{S\'eminaire Bourbaki}, exp.~221 (1960/61),
collected edition, vol.~6, pp.~249--276.
Theorem 3.2, p.~260; Proposition 3.8, p.~265.
\url{https://www.numdam.org/item/SB_1960-1961__6__249_0/}.

\bibitem{Hironaka}
H. Hironaka,
Resolution of singularities of an algebraic variety over a field of
characteristic zero, I, II,
\emph{Ann. of Math.} (2) \textbf{79} (1964), 109--203, 205--326.
Part I, Main Theorems I and I$^*$, pp.~132--133.
\url{https://doi.org/10.2307/1970486};
\url{https://doi.org/10.2307/1970547}.

\bibitem{Huybrechts}
D. Huybrechts,
\emph{Lectures on K3 Surfaces},
Cambridge Studies in Advanced Mathematics 158,
Cambridge University Press, 2016.
Chapter and result locators in this paper refer to the author's
prepublication draft:
\url{https://www.math.uni-bonn.de/people/huybrech/K3Global.pdf}.

\bibitem{Ehresmann}
I. Kol\'a\v r, J. Slov\'ak, and P.~W. Michor,
\emph{Natural Operations in Differential Geometry},
Springer, 1993, Lemma 9.2, p.~77.
Corrected electronic edition:
\url{https://www.mat.univie.ac.at/~michor/kmsbookh.pdf}.

\bibitem{KugaSatake}
M. Kuga and I. Satake,
Abelian varieties attached to polarized K3-surfaces,
\emph{Math. Ann.} \textbf{169} (1967), 239--242.
\url{https://doi.org/10.1007/BF01399540}.

\bibitem{Milne}
J.~S. Milne,
\emph{Introduction to Shimura Varieties},
revised electronic version, 16 September 2017.
Theorems 6.7--6.8 and Corollary 6.9, pp.~72--73.
\url{https://www.jmilne.org/math/xnotes/svi.pdf}.

\bibitem{SerreGAGA}
J.-P. Serre,
G\'eom\'etrie alg\'ebrique et g\'eom\'etrie analytique,
\emph{Ann. Inst. Fourier} \textbf{6} (1956), 1--42.
\S12, Theorems 1--3, pp.~19--20; \S19, Propositions 13--15,
pp.~29--30.
\url{https://doi.org/10.5802/aif.59}.

\bibitem{StacksCartier}
The Stacks Project Authors,
\emph{The Stacks Project},
Lemma 42.29.5, tag 02TB (Cartier intersection and fundamental cycles).
\url{https://stacks.math.columbia.edu/tag/02TB}.

\bibitem{StacksCorrespondences}
The Stacks Project Authors,
\emph{The Stacks Project},
\S45.3, tag 0FFZ (correspondences), especially
Lemma 45.3.1 and Lemma 45.3.6.
\url{https://stacks.math.columbia.edu/tag/0FFZ}.

\bibitem{StacksFamilies}
The Stacks Project Authors,
\emph{The Stacks Project},
tags 01V4 and 01V8 (smoothness and the fiber criterion),
01VU (high powers of ample sheaves),
and 0B5U (ampleness and cohomology).
\url{https://stacks.math.columbia.edu/tag/01V4};
\url{https://stacks.math.columbia.edu/tag/01V8};
\url{https://stacks.math.columbia.edu/tag/01VU};
\url{https://stacks.math.columbia.edu/tag/0B5U}.

\bibitem{Voisin}
C. Voisin,
\emph{Hodge Theory and Complex Algebraic Geometry I},
translated by L. Schneps,
Cambridge Studies in Advanced Mathematics 76,
Cambridge University Press, 2002.
\S\S6.1--6.3 and 7.1--7.3.
\url{https://doi.org/10.1017/CBO9780511615344}.

\end{thebibliography}

\end{document}
